\documentclass{article}

\usepackage{enumitem}

\author{Franz Scharnreitner, Jakob Kapeller}
\title{The effects of a speculative assset price on the supply of money in a post keynsian framework}
\begin{document}

\maketitle

\section{Introduction}

\begin{enumerate}[label=(\alph*)]
\item model provides a missing link between debates on financialization ond PK models that are mostly focused on the real economy.
\item the notion of specullation and liquidty preference has always been there, but whether and to what extent it has an impact is debate. while for mainstreamers money as a whole is only a vial, Keynes emphaizses that the demand for money –  e.g. in the form of hoarded savings – can have real impacts. however, in treatise of money he denies that the supply of money for speculative purposes has such impact. our approach shows what happens if there is such a link and clarifies the most obvious assumptions that come with imposing one in a general framework of endogenous money.
\item our argument on turnover implicitly creates a multiplier that governs the mismatch in scale between speculative transactions and GDP (also already discussed in the Treatise!)
\end{enumerate}
\section{A simple static model}

\subsection{The core model}


Three versions:

\begin{enumerate}[label=(\alph*)]
  \item  In the first one just speculation is added and it is shown how this inflates balance sheets.
  \item  In the second one we impoase the link via c to show this has real impact. we could compare this with a scenario of rising inequality.
  \item  we do the same as in c with a different notion of causality,
where increasing totals bank balance sheets constrain lending to that the reaction of the banking sector is modeled in greater detaill. should give a result similar to (b) that is better explained.
\end{enumerate}


\subsection{Endogenizing $\alpha$ - towards an endogenous asset turnover rate}

\subsection{Endogenizing $c$ - credit constrains through speculation}
\subsubsection{Dual interest rate policies}

\subsection{Policies}
→ for this we need some 2/4 standard panel to Xplain results.

\section{Steps towards a dynamic model}



\end{document}
